\chapter{序論}

\section{本章の概要}
\begin{figure}[tbp]
\centering
\includegraphics[width=0.92\linewidth]{framework_overview}
\caption{本研究の全体像。特許文書を課題文脈と解決手段文脈に分離し、専門家確認候補を抽出するスクリーニング手法として整理する。}
\label{fig:framework_overview}
\end{figure}

本章では、本研究の背景、問題意識、研究目的、本研究で取り組む課題、対象範囲、貢献、および本論文の構成を述べる。本研究のテーマは、特許文書の課題・解決手段分離に基づくクロスドメイン技術候補抽出である。

特許文書には、発明が解決しようとする課題、具体的な解決手段、対象物、効果などが記述されている。そのため、特許文書は技術動向を把握するための情報源であるだけでなく、異分野に存在する技術手法や解決手段を探索するための知識源としても利用できる。

一方で、特許文書全体を一つの表現として扱う場合、課題、対象物、解決手段、効果が混在し、候補が何に基づいて抽出されたのかを説明しにくい。特に、課題が近い異分野間において、対象分野では低出現または未使用である技術手法・解決手段を探索するには、課題文脈と解決手段文脈を分けて扱う必要がある。

そこで本研究では、特許文書を課題文脈である \texttt{problem\_context} と解決手段文脈である \texttt{solution\_context} に分離し、課題が近い異分野間において専門家が確認すべき技術候補を抽出する枠組みを検討する。本研究は、技術導入可能性や将来の技術移転を証明するものではなく、特許文書に基づいて専門家確認候補を絞り込むためのスクリーニング手法として位置づけられる。

\section{研究背景}

本節では、本研究の背景として、特許文書を用いた異分野間の技術候補探索の意義を述べる。

企業や研究機関が新しい技術開発の方向性を検討する際、自分たちの分野内の技術だけでなく、他分野で用いられている技術手法や解決手段を参照することが重要になる場合がある。ある分野で発展した技術手法が、別分野の類似課題に対する解決手段となる可能性があるためである。たとえば、画像解析、材料処理、洗浄、検査、制御などの技術では、対象物や用途が異なっていても、領域抽出、異常検出、残渣除去、表面処理、精密制御といった共通する課題が含まれることがある。

このような異分野間の技術候補探索において、特許文書は有用な情報源である。特許文書には、技術の名称だけでなく、発明が何を目的としているのか、どのような課題を解決しようとしているのか、どのような手段を用いているのか、どのような効果が期待されるのかが記述されている。したがって、特許文書を分析することで、単なる技術語の出現頻度だけではなく、課題と解決手段の関係を読み取ることができる。

近年では、自然言語処理、文書埋め込み技術、大規模言語モデルの発展により、大量の特許文書を用いた分析の範囲が広がっている。特許文書をベクトル化し、意味的類似度に基づいて検索・ランキングすることで、従来よりも広い範囲から類似特許や関連技術を探索できるようになった。また、特許文書から課題、解決手段、対象物、効果などを抽出し、技術動向分析や技術機会探索に用いる試みも進んでいる。

一方で、特許文書は長く、複数の情報が混在している。発明の課題、対象物、構成、作用効果、請求項の限定表現が一つの文書に含まれるため、文書全体を一つの表現として扱った場合、どの情報が類似度に強く影響したのかを説明しにくい。この点は、課題が近い異分野間において、対象分野では低出現または未使用である技術手法・解決手段を探索する際に特に問題となる。そこで次節では、本研究が対象とする問題意識として、全文類似に基づく候補抽出の限界と、課題文脈と解決手段文脈を分けて扱う必要性を整理する。

\section{問題意識}

本節では、本研究が対象とする問題意識を整理する。

異分野間の技術候補探索では、対象分野と参照分野で用いられる語彙、対象物、用途名が異なることが多い。そのため、キーワード一致だけでは、同じような課題に取り組んでいる候補を見落とす可能性がある。一方で、全文類似を用いる場合でも、特許文書全体の類似度が高い理由が、課題の近さによるものなのか、対象物や用途名の一致によるものなのか、解決手段の類似によるものなのかを区別しにくい。

この問題は、課題が近い異分野間において、対象分野では低出現または未使用である技術手法・解決手段を探索する際に特に重要である。ある候補特許が対象分野の特許と全文として類似していたとしても、その類似が対象物や分野名の一致に由来するだけであれば、異分野の解決手段を確認する候補としては十分でない場合がある。逆に、対象物や用途名は異なっていても、解いている課題が近く、参照分野で用いられている解決手段が対象分野の課題に対して参考になりうる候補は、全文類似だけでは上位に現れにくい可能性がある。

また、候補抽出結果を人間が確認する段階では、候補が提示された理由を説明できることが重要である。評価者や専門家は、単にスコアが高い候補を見るだけでなく、「どの課題が近いのか」「どの解決手段が候補として提示されているのか」を理解する必要がある。全文類似に基づくランキングでは、課題、対象物、解決手段、効果が混在しやすいため、候補を専門家確認候補として提示する際の根拠を整理しにくい。

したがって、クロスドメイン技術候補抽出では、特許文書全体を一つの表現として扱うだけでなく、課題文脈と解決手段文脈を分けて扱う必要がある。課題文脈によって「何の課題が近いのか」を捉え、解決手段文脈によって「どのような解決手段が用いられているのか」を確認することで、対象分野では低出現または未使用である技術手法・解決手段を、専門家が確認すべき候補として説明しやすく提示できると考えられる。

\section{研究目的}

本節では、本研究の目的を述べる。

本研究の目的は、特許文書を課題文脈と解決手段文脈に分離し、課題が近い異分野間において、対象分野では低出現または未使用である技術手法・解決手段を、専門家確認候補として抽出する枠組みを設計し、実データを用いて確認することである。

本研究では、特許文書の Title、Abstract、Claims を中心に用い、LLM によって \texttt{problem\_context}、\texttt{solution\_context}、\texttt{technical\_means}、\texttt{target\_object}、\texttt{application\_field}、\texttt{effect\_context} などの項目を抽出する。ここで、\texttt{problem\_context} は発明が解決しようとする課題を表し、\texttt{solution\_context} はその課題に対して提示される解決手段を表す。LLM は候補の妥当性を最終判定するものではなく、長く複雑な特許文書を、人間が確認しやすい文脈単位に整理するための支援として用いる。

抽出された文脈を用いて、参照分野と対象分野の間で課題文脈の類似性を確認し、解決手段文脈および \texttt{technical\_means} に基づいて、対象分野では低出現または未使用である技術手法・解決手段を候補としてランキングする。この手順により、単に全文が似ている特許を抽出するのではなく、どの課題が近く、どの解決手段が候補として提示されるのかを説明しやすくする。

本研究では、この枠組みを二つの実験により確認する。実験1では、医療画像解析分野で先行して用いられた U-Net を対象に、製造欠陥検出分野への後年出現を参照した後ろ向き評価を行う。実験1は、本研究における近接解決策型候補抽出の主分析である。実験2では、洗浄・除去系技術から半導体洗浄への異質解決策型候補抽出を補助的に分析する。実験2は、対象物や工程が大きく異なる場合にも、課題文脈を手がかりに専門家確認候補を整理できるかを確認するための補助分析である。

なお、本研究で抽出する候補は、技術導入可能性や将来の技術移転が確認されたものではない。本研究は、特許文書に基づいて専門家が確認すべき候補を絞り込むためのスクリーニング手法として位置づける。

\section{本研究で取り組む課題}

本研究の目的を達成するために、本研究では以下の3点を課題として設定する。

第1に、クロスドメイン技術候補抽出において、特許文書から課題文脈と解決手段文脈を分けて扱う枠組みを定義する必要がある。特許文書には、発明が解決しようとする課題、具体的な解決手段、対象物、効果などが含まれるため、文書全体を一つの表現として扱うだけでは、候補抽出の根拠を明確に説明しにくい。そのため、課題文脈と解決手段文脈を区別し、異分野間の技術候補探索に用いるための分析枠組みを整理する必要がある。

第2に、特許文書から課題文脈と解決手段文脈を抽出し、課題文脈の類似性と解決手段文脈の差分に基づいて候補をランキングする方法を設計する必要がある。課題が近い異分野間において、参照分野では用いられているが対象分野では低出現または未使用である技術手法・解決手段を抽出するためには、単なる全文類似ではなく、課題の近さと解決手段の違いを分けて扱う必要がある。

第3に、提案する枠組みが実際の特許データにおいて専門家確認候補を抽出できるかを、具体的な事例を用いて確認する必要がある。本研究では、医療画像解析分野から製造欠陥検出分野への U-Net の後年出現を参照した実験と、洗浄・除去系技術から半導体洗浄への異質解決策型候補抽出を対象とした補助的実験を行い、提案手法により抽出された候補の妥当性を確認する。

\section{本研究の対象範囲}

本節では、本研究で扱う範囲と扱わない範囲を明確にする。

本研究が扱う対象は、特許文書に記述された課題、解決手段、対象物、効果などの情報を用いた、クロスドメイン技術候補の抽出である。ここでいう候補とは、対象分野において専門家が確認すべき技術手法・解決手段の候補であり、実際に導入可能であることが確認された技術を意味するものではない。

本研究は、特許文書に基づいて専門家確認候補を絞り込むためのスクリーニング手法を対象とする。そのため、実験結果は、候補の解釈可能性や確認対象としての妥当性を検討するものであり、技術導入可能性、製品化可能性、事業化可能性を証明するものではない。また、将来の技術移転や企業進出を予測することも目的としない。特許文書は技術知識源として有用である一方で、特許が存在することは、その技術が実装されたことや製品化されたことを直接意味しないためである。

評価については、LLM支援付き著者確認評価として扱う。これは、LLMによって整理された課題文脈、解決手段文脈、技術手段、対象物、効果などの情報をもとに、著者が候補を確認する評価である。したがって、本研究の評価は、完全な第三者専門家評価ではない。この点は、本研究の限界として明示する。

また、本研究では、近接解決策型と異質解決策型という二つの候補タイプを扱う。近接解決策型は、参照分野と対象分野の課題や解決手段が比較的近く、ある分野で先行して現れた手法が別分野でも後年確認されるような候補を想定する。異質解決策型は、対象物、用途、材料、工程が大きく異なる一方で、課題や効果の観点から専門家確認候補として読みうる技術を想定する。本研究では、実験1を近接解決策型の主分析、実験2を異質解決策型の補助分析として位置づける。

\section{本研究の貢献}

本節では、本研究の貢献を整理する。

第一の貢献は、特許文書を課題文脈である \texttt{problem\_context} と解決手段文脈である \texttt{solution\_context} に分離し、クロスドメイン技術候補を抽出する枠組みを示した点である。全文類似に基づく候補抽出では、課題、対象物、解決手段、効果が混在しやすく、候補が何に基づいて抽出されたのかを説明しにくい。これに対して本研究では、課題文脈によって「何の課題が近いのか」を捉え、解決手段文脈や \texttt{technical\_means} によって「どの解決手段が候補として提示されているのか」を確認できる形に整理する。これにより、対象分野では低出現または未使用である技術手法・解決手段を、専門家確認候補として説明しやすく提示する枠組みを示した。

第二の貢献は、U-Net を対象とした近接解決策型の主分析を行った点である。医療画像解析分野で先行して用いられた U-Net を対象に、製造欠陥検出分野への後年出現を参照した後ろ向き評価を行い、過去期間の特許文書に基づいて U-Net 関連技術を専門家確認候補として抽出できるかを確認した。この分析により、課題と解決手段が比較的近い異分野間において、提案する枠組みが候補抽出の根拠を説明しやすい形で提示できる可能性を示した。ただし、この結果は医療画像解析分野から製造欠陥検出分野への技術移転を証明するものではない。

第三の貢献は、異質解決策型の補助分析と、LLM支援付き著者確認評価に基づく評価設計を行った点である。洗浄・除去系技術から半導体洗浄への候補抽出を補助的に分析し、対象物、用途、材料、工程が大きく異なる場合にも、課題文脈を手がかりとして専門家確認候補を整理できるかを確認した。また、抽出された課題文脈、解決手段文脈、技術手段、対象物、効果などを確認できる形で評価用データを作成し、LLM支援付き著者確認評価を行った。これにより、候補の良し悪しを自動的に断定するのではなく、人間が確認すべき情報を整理するスクリーニング手法としての評価設計を示した。

以上の貢献により、本研究は、特許文書を用いた異分野技術候補探索において、全文類似のみでは説明しにくい候補抽出理由を、課題文脈と解決手段文脈に分けて整理する枠組みを示した。本研究の貢献は、技術導入可能性や将来の技術移転を証明することではなく、専門家が確認すべきクロスドメイン技術候補を説明可能な形で絞り込む点にある。

\section{本論文の構成}

本節では、本論文の構成を述べる。

第1章では、本研究の背景、問題意識、研究目的、本研究で取り組む課題、対象範囲、貢献、および論文構成を述べた。特許文書を用いたクロスドメイン技術候補抽出において、全文類似だけでは課題の近さと解決手段の差分を分けて捉えにくいという問題意識を示し、本研究を専門家確認候補を絞り込むためのスクリーニング手法として位置づけた。

第2章では、本研究に関連する研究を整理する。具体的には、特許文書を用いた技術分析、特許文書埋め込みと類似度分析、技術機会分析、類推型特許マイニング、技術拡散・時間差分析、LLMを用いた技術文書の構造化に関する研究を概観し、本研究の位置づけを明確にする。

第3章では、提案手法を説明する。特許文書から課題文脈である \texttt{problem\_context}、解決手段文脈である \texttt{solution\_context}、\texttt{technical\_means} などを抽出し、課題文脈の類似性と解決手段・技術手段の差分に基づいて候補をランキングする流れを述べる。また、本研究で扱う近接解決策型と異質解決策型の候補タイプを定義する。

第4章では、データセットと実験設計を説明する。実験1で用いる医療画像解析分野および製造欠陥検出分野のデータ、実験2で用いる洗浄・除去系技術および半導体洗浄関連データを整理し、期間分割、ベースライン比較、評価方法を述べる。

第5章では、実験1として、U-Netを対象とした近接解決策型候補抽出を検証する。医療画像解析分野で先行して用いられた U-Net を対象に、製造欠陥検出分野への後年出現を参照した後ろ向き評価を行い、提案手法が専門家確認候補を抽出できるかを確認する。実験1は本研究の主分析である。

第6章では、実験2として、洗浄・除去系技術を対象とした異質解決策型の補助分析を行う。洗浄・除去系技術から半導体洗浄への候補抽出を対象に、対象物や工程が大きく異なる場合にも、課題文脈を手がかりとして専門家確認候補を整理できるかを確認する。実験2は、本研究の補助分析である。

第7章では、実験1と実験2の結果を踏まえて考察する。課題文脈と解決手段文脈を分離することの有効性、全文類似に基づく候補抽出との違い、スクリーニング手法としての意義、本研究の限界、および今後の展望を整理する。

第8章では、本研究の結論を述べる。研究全体をまとめ、得られた示唆、限界、今後の課題を整理し、特許文書に基づくクロスドメイン技術候補抽出における課題・解決手段分離の位置づけをまとめる。

第9章では、本研究の理解を補助するための参考図表を示す。また、付録では、LLMによる文脈抽出プロンプトの概要を示す。

以上の構成により、本論文では、特許文書を全文として一括で扱うのではなく、課題文脈と解決手段文脈に分離することで、異分野間の専門家確認候補を説明しやすく抽出する枠組みを検討する。
